% !TeX root = ../../main.tex
\subsection{道のり・速さ・時間の問題}

    基本は，道のり・速さ・時間の関係の考え方を使って，\uwave{進んだ道のりは変わらない}ことから，

    \[\textbf{(道のりの合計) = (道のりの合計)}\]

    や，\uwave{かかった時間は変わらない}ことから，

    \[\textbf{(かかった時間の合計) = (かかった時間の合計)}\]

    \vspace{1.5em}
    という式を立てる.特に，速さ・時間・距離の問題のうち，代表的なものに
    \uwave{「追いつく」問題}がある.
    \vspace{1em}

    \begin{techniquebox}[tech:03]{\textbf{「追いつく」問題}}
        \textbf{(Aが進んだ道のり) = (Bが進んだ道のり)} という考え方を使う!
    \end{techniquebox}

    \vspace{1em}
    Technique \ref{tech:03} を応用した，\uwave{「池の周りをまわる問題」}もある.

    \begin{example}[【例題 2】]
        家から1500mのところにある学校に向かいました.分速50mで歩いていましたが，
        途中から分速100mで走りました.合計で20分で着きました.
        歩いた道のりは何mか求めなさい.
    \end{example}

    \begin{proof}\mbox{}\\
        \[（歩いた時間）+（走った時間）=（合計でかかった時間）\]

        という方針にしたがって，方程式を立てる.

        \begin{hiddenanswer}
            \begin{align*}
                \text{歩いた道のりを $x$ mとすると，}\\
                \dfrac{x}{50} + \dfrac{1500 - x}{100} &= 20 \\
                2x + 1500 - x &= 2000 \\
                x &= 500
            \end{align*}

            \finalanswer{答え：500m}
        \end{hiddenanswer}
    \end{proof}
\newpage
